\documentclass{article}
\usepackage{tcee-repaso}
\tcetema[Perspectivas y flujos de la economía mundial]{3.B.32}{Perspectivas económicas mundiales. Estructura sectorial y geográfica de los flujos comerciales y financieros internacionales. Las nuevas áreas emergentes.}
\tcefecha{14/11/2026}
\begin{document}
\cabecerarepaso
\begin{columnas}

% ==========================================================================
\bloque{Esquema del cante}
\esquemacante{%
\apartadocante{Introducción: «Dios aprieta, pero no ahoga»}{2,5}
\apartadocante{I. Estructura de los flujos}{13}
\subcante{1. Sistema comercial: mercancías, servicios, gravedad}{7}
\subcante{2. Sistema financiero: capital, factores, ciclo global}{4,5}
\subcante{3. Sistema monetario: el dólar}{1,5}
\apartadocante{II. Las nuevas áreas emergentes}{5}
\subcante{1. Concepto y causas: convergencia condicional}{2,5}
\subcante{2. Principales áreas}{1,5}
\subcante{3. Retos}{1}
\apartadocante{III. Perspectivas económicas mundiales}{8,5}
\subcante{1. Coyuntura: resistente bajo presión}{2}
\subcante{2. Previsiones de los organismos}{1,5}
\subcante{3. Riesgos y recomendaciones}{1,5}
\subcante{4. Medio y largo plazo: centro de gravedad}{2,5}
\subcante{5. Balance: DAFO}{1}
\apartadocante{Conclusión}{1}}
\emph{Gráficos del cante}: recta de gravedad (\hyperref[rep:I.1]{I.1}), circuito del ciclo financiero global (\hyperref[rep:I.2]{I.2}), nube de convergencia (\hyperref[rep:II.1]{II.1}), centro de gravedad y DAFO (\hyperref[rep:III.4]{III.4}-\hyperref[rep:III.5]{III.5}). Funciones clave: la ecuación de gravedad y la de convergencia $\beta$. Una cifra por idea.

% ==========================================================================
\bloque[2,5]{Introducción}
\begin{itemize}
\item \textbf{Enganche}: «Dios aprieta, pero no ahoga»: en 2026, guerra en Oriente Medio, Ormuz cerrado (1/4 del crudo marítimo), petróleo +41\,\%, aranceles de EE.\,UU. máximos desde los 40… y el PIB mundial crece \dato{3,1\,\%}{FMI, WEO oct. 2026}, con la IA y Asia. \emph{¿Hacia dónde va la economía mundial y quién la mueve?}
\item \textbf{Relevancia}: contexto de toda decisión de política; España, abierta: exportaciones, turismo e IED dependen de la demanda mundial \vertema{4.A.22} \vertema{4.A.23} \vertema{4.A.29}; el oficio del técnico comercial.
\item \textbf{Contextualización}: \concepto{flujos comerciales} (bienes y servicios), \concepto{financieros} (IED, cartera, otra inversión, reservas) \vertema{3.B.11} y \concepto{de factores} (trabajo, tecnología); estructura sectorial (qué) y geográfica (quién con quién); \emph{nuevas áreas emergentes}; \emph{perspectivas} (FMI, OCDE, Comisión, BM, OMC); tras la hiperglobalización, lentiglobalización y fragmentación \vertema{3.B.21}.
\item \textbf{Problemática}: ¿Qué se comercia y por qué importa la distancia? ¿Cómo se mueven el capital y los factores, y en qué moneda? ¿Quiénes emergen y por qué? ¿Cómo está, qué prevén los organismos y qué riesgos?
\item \textbf{Estructura}: del presente al futuro: \hyperref[rep:I]{I} flujos (comercial, financiero, monetario, como en el sistema económico internacional \vertema{3.B.20} \vertema{3.B.21}); \hyperref[rep:II]{II} emergentes; \hyperref[rep:III]{III} perspectivas.
\end{itemize}

% ==========================================================================
\bloque[13]{I. Estructura de los flujos\label{rep:I}}
\begin{ideaclave}
Comercio: manufacturas, cada vez más servicios y bienes de la IA; 3 polos con Asia en ascenso; la distancia pesa (gravedad). Finanzas: volátiles, entre ricos, al ritmo del ciclo financiero global. Personas y tecnología, mucho menos libres. Todo, sobre todo en dólares.
\end{ideaclave}

\sub{I.1. Sistema comercial\label{rep:I.1}}
\emph{Mercancías}
\begin{itemize}
\item \dato{26,3 billones \$ (+7\,\%)}{OMC, 2025}; servicios 9,6 (+8\,\%); total $\approx$35 billones, $\approx$30\,\% del PIB mundial.
\item \emph{Etapas}: hiperglobalización (1991-2007: comercio 6,2\,\% frente a PIB 3,3\,\%); colapso de 2009 ($>$10\,\%) y lentiglobalización (3,2 frente a 3,0); choques y resistencia.
\item \emph{CUCI}: manufacturas \dato{69,7\,\%}{BM, 2025} (75 en 1995); combustibles 9,5, minerales 4 (suben con los precios); alimentos 9, agrícolas 1.
\item \emph{Geografía}: Asia oriental 22 $\to$ 35\,\%; Europa y Asia central 51 $\to$ 38\,\%; América del Norte 15 $\to$ 10\,\%. China 1.\textsuperscript{er} exportador desde 2009 (\dato{14,4\,\%}{OMC}); EE.\,UU. 1.\textsuperscript{er} importador (13,2\,\%); España 20.ª (1,7\,\%). Asia, 71\,\% del crecimiento de 2025. Acuerdos regionales: mercado único, T-MEC, RCEP, ZLCCA.
\item \emph{Formas}: interindustrial (ventaja comparativa); \concepto{intraindustrial} (países ricos y parecidos; índice de \autor{Grubel}-\autor{Lloyd}, $1-\lvert X_k-M_k\rvert/(X_k+M_k)$); \concepto{intraempresa} (filiales, precios de transferencia); bienes intermedios (CGV).
\end{itemize}
\graficorepaso[0.9\linewidth]{comercio-productos}{Exportaciones de mercancías por productos (CUCI), 1995-2025}
\graficorepaso[0.9\linewidth]{exportaciones-regiones}{Cuota en las exportaciones mundiales de mercancías, 1990-2025}

\emph{Servicios y cadenas globales de valor}
\begin{itemize}
\item 1/4 del comercio (casi 1/2 en valor añadido: servitización); 4 modos del AGCS (el 3, presencia comercial, el principal) \vertema{3.B.34}.
\item Otros servicios comerciales 60\,\% (informáticos 13\,\%), viajes 20\,\%, transporte 16\,\%; \concepto{servicios prestados digitalmente} 5,3 billones (15\,\% del comercio). EE.\,UU. 12,7\,\%, India 4,3\,\%, España 10.ª (2,6\,\%) \vertema{4.A.13}.
\item \concepto{CGV}: fases en $\geq 2$ países; comercio en valor añadido (TiVA); 3 «fábricas»; se reordenan (\emph{friend-shoring}, \emph{near-shoring}) \vertema{3.B.18}.
\end{itemize}
\graficorepaso[0.9\linewidth]{servicios-composicion}{Servicios comerciales por tipo, 2005-2025: TIC 7,1 $\to$ 15,7\,\%}

\emph{Geografía: el modelo de gravedad}
\begin{itemize}
\item \emph{a) Idea}: \autor{Tinbergen} (1962; Nobel 1969): comercio $\uparrow$ con el tamaño, $\downarrow$ con la distancia \vertema{3.B.10}.
\item \emph{b) i) Supuestos}: gasto proporcional a la renta; bienes diferenciados por origen; comerciar cuesta con la distancia. \emph{ii) Desarrollo}:
\begin{formula}
\centering \funcionclave{\ln X_{ij}=\beta_0+\beta_1\ln Y_i+\beta_2\ln Y_j-\beta_3\ln d_{ij}+\beta_4\,\mathrm{ACR}_{ij}+u_{ij}}
\end{formula}
\item[] \autor{Anderson} (1979), \autor{Anderson}-\autor{van Wincoop} (2003): resistencia multilateral (efectos fijos). \emph{iii) Implicaciones}: elasticidad renta $\approx1$, distancia $\approx-1$; fronteras (\autor{McCallum}: $\times$20) y acuerdos.
\item \emph{c) Extensiones}: \autor{Krugman} (1980; Nobel 2008), \autor{Melitz}, \autor{Chaney} (margen extensivo); la geopolítica como distancia (\autor{Gopinath} et al.).
\item \emph{d) Evidencia}: España 2024, 160 destinos: pendiente $-1{,}29$ ($R^2=0{,}39$), UE por encima; la distancia no pierde peso (\autor{Disdier}-\autor{Head}, 2008).
\item \emph{e) Valoración}: ajusta muy bien y evalúa la política comercial; la distancia mezcla costes y dice cuánto, no qué.
\end{itemize}
\graficorepaso[0.92\linewidth]{gravedad-espana}{Gravedad: exportaciones de España por destino y distancia, 2024}

\sub{I.2. Sistema financiero\label{rep:I.2}}
\emph{Flujos financieros}
\begin{itemize}
\item Cuenta financiera: IED ($\geq$10\,\%, estable), cartera (volátil), otra inversión (préstamos, procíclica), derivados, reservas.
\item Entre ricos y centros financieros; flujos brutos $\approx$20\,\% del PIB mundial (2007) y desplome; capital «cuesta arriba» (\concepto{paradoja de \autor{Lucas}}, 1990; Nobel 1995) \vertema{3.B.12}.
\item \emph{IED}: \dato{1,6 billones \$ (+6\,\%)}{UNCTAD, 2025}; en desarrollo 0,9 (Asia 40\,\% del mundo); sectores estratégicos 44\,\% del \emph{greenfield} (16\,\% en 2020), fusiones y adquisiciones $-7$\,\%; los pobres, $\approx$10\,\%. Por qué: paradigma OLI (\autor{Dunning}, 1980).
\item \emph{Banca}: \dato{49,5 billones \$}{BPI, 2026T2} (32,9 en 2007); crédito +11\,\%; emergentes $\approx$10\,\%.
\item \emph{Desequilibrios}: saldos 3,8\,\% PIB mundial; China +735.000 M\$; EE.\,UU. $-1{,}1$ billones; PIIN EE.\,UU. $-27{,}5$ billones \vertema{4.A.21}.
\end{itemize}
\graficorepaso[0.9\linewidth]{ied-mundial}{IED: entradas netas por grupos de renta, 1990-2025 (BM)}
\graficorepaso[0.9\linewidth]{bancos-transfronterizos}{Activos transfronterizos de los bancos, 1999-2026 (BPI)}

\emph{Ciclo financiero global}
\begin{itemize}
\item \emph{a) Idea}: \autor{Rey} (2013): precios de activos, flujos brutos y crédito se mueven a la vez en todo el mundo; lo dirigen la Fed y la aversión al riesgo.
\item \emph{b) i) Supuestos}: capital libre e intermediarios globales en dólares; su restricción depende de la volatilidad; dólar dominante. \emph{ii) Desarrollo}: 1) la Fed relaja (VIX $\downarrow$); 2) se apalancan; 3) entran capitales, crédito y precios $\uparrow$; 4) la Fed endurece: salidas y paradas súbitas. \emph{iii) Implicaciones}: \emph{dilema, no trilema} (\autor{Mundell}, Nobel 1999): flotar no aísla; controles/macroprudencial o menos autonomía \vertema{3.B.15}.
\item \emph{c) Extensiones}: \autor{Miranda-Agrippino}-\autor{Rey} (2020); moneda dominante (\autor{Gopinath}).
\item \emph{d) Evidencia}: 2009-2012, \emph{taper tantrum} (2013), marzo de 2020, 2022; en 2026, crédito +11\,\% pero Asia emergente $-2{,}9$\,\%.
\item \emph{e) Valoración}: explica la sincronía; pero los que flotan con buenas instituciones conservan autonomía (\autor{Obstfeld} et al., 2019).
\end{itemize}
\graficorepaso[0.85\linewidth]{ciclo-global}{El ciclo financiero global}

\emph{Flujos de factores}
\begin{itemize}
\item \emph{Trabajo}: \dato{304 millones de migrantes}{ONU, 2024}, el 3,7\,\% de la población (2,9\,\% en 1990); remesas 685.000 M\$ (2024) $>$ IED a renta baja y media; modo 4, el más restringido.
\item \emph{Tecnología}: incorporada en bienes e IED; cobros por propiedad intelectual \dato{629.000 M\$}{BM, 2025} (402.000 en 2018).
\end{itemize}

\sub{I.3. Sistema monetario\label{rep:I.3}}
\begin{itemize}
\item \emph{Dólar dominante}: \dato{56,7\,\%}{COFER, jun. 2026} de las reservas asignadas (71\,\% en 1999); euro 20,6\,\%, renminbi 2,1\,\%; 89,2\,\% de las operaciones de divisas; moneda dominante en la facturación \vertema{3.B.13}.
\item \emph{Diversificación lenta}: oro ($>$1.000 t/año desde 2022; 27\,\% de las reservas a precios de mercado), monedas «no tradicionales». Renminbi sin convertibilidad; \emph{stablecoins} en dólares (ley GENIUS, 2025) \vertema{3.B.21}.
\end{itemize}

% ==========================================================================
\bloque[5]{II. Las nuevas áreas emergentes\label{rep:II}}
\begin{ideaclave}
Crecen más porque parten de más lejos (convergencia condicional), pero solo convergen los que elevan su estado estacionario; Asia lidera, el resto sigue atado a las materias primas.
\end{ideaclave}

\sub{II.1. Concepto y causas\label{rep:II.1}}
\begin{itemize}
\item Sin definición oficial (CFI, años 80). FMI: avanzadas/EMED; BM: renta baja ($\leq$1.175\,\$), media-baja, media-alta ($\leq$14.375\,\$), alta; MSCI/FTSE: mercados. \textbf{BRICS}: \autor{O'Neill} (2001); +Sudáfrica (2010), Egipto, Etiopía, Irán, EAU (2024), Indonesia (2025), Arabia Saudí invitada; 10 socios; Nueva Delhi (sep. 2026): \dato{40,5\,\% del PIB en PPA}{FMI, 2025}.
\item \emph{Peso} (PPA): avanzadas 65\,\% (1980) $\to$ 39,5\,\% (2025); EMED por delante desde 2009; China 19,6\,\% (supera a EE.\,UU. en 2016), India 8,2 $\to$ 10,2\,\% (2031). A tipos de mercado: EE.\,UU. 26\,\%, China 16,6\,\%.
\item \emph{a) Idea}: \autor{Solow}: crece más quien está lejos de su estado estacionario, que depende de ahorro, población, capital humano e instituciones \vertema{3.A.45}.
\item \emph{b) i) Supuestos}: tecnología común, rendimientos decrecientes; estados estacionarios distintos. \emph{ii) Desarrollo}:
\begin{formula}
\centering \funcionclave{\tfrac{1}{T}\ln\tfrac{(Y/L)_{i,t}}{(Y/L)_{i,t-T}}=\beta_0+\beta_1\ln\left(\tfrac{Y}{L}\right)_{i,t-T}+\vect{x}_i\cdot\vect{\psi}+u_{i,t}}
\end{formula}
\item[] absoluta ($\beta_1<0$ con $\vect{\psi}=\vect{0}$) o condicional (\autor{Barro}-\autor{Sala-i-Martín}; MRW, 1992). \emph{iii) Implicaciones}: emergen los que elevan su estado estacionario; se frena al acercarse.
\item \emph{c) Evidencia}: desde 2000, convergencia absoluta lenta: 189 países, pendiente $-0{,}43$, $R^2=0{,}08$; China 7,7\,\%, Vietnam 5,1\,\%, India 4,9\,\% anual; instituciones (\autor{Acemoglu}-\autor{Johnson}-\autor{Robinson}, Nobel 2024).
\item \emph{d) Valoración}: explica quién emerge; pero causalidad dudosa y falacia de \autor{Galton}.
\end{itemize}
\graficorepaso[0.9\linewidth]{peso-pib-ppa}{Peso en el PIB mundial en PPA, 1980-2031}
\graficorepaso[0.92\linewidth]{convergencia}{Convergencia: PIB por habitante (2000) y crecimiento 2000-2024}

\sub{II.2. Principales áreas\label{rep:II.2}}
\begin{itemize}
\item \textbf{China}: 19,6\,\% PIB (PPA), 14,4\,\% de las exportaciones; 4,5 $\to$ 4\,\%; inmobiliario, deuda local, superávit récord; vehículos eléctricos, baterías, renovables.
\item \textbf{India}: 7\,\%; 8,2 $\to$ 10,2\,\% del PIB mundial; servicios informáticos; electrónica.
\item \textbf{Sudeste Asiático}: ASEAN-5 4,4\,\%, Vietnam $>$8\,\%: economías conectoras \vertema{3.B.35}.
\item \textbf{Golfo}: fondos soberanos y diversificación; Arabia Saudí $-3{,}3$\,\% (2026), rebote en 2027.
\item \textbf{América Latina}: 2,3\,\%; materias primas, minerales críticos, \emph{near-shoring} (México); UE-Mercosur \vertema{3.B.41}.
\item \textbf{África subsahariana}: 4,5\,\%, la población más joven; ZLCCA ($\approx$50 ratificaciones).
\item Sur-Sur: +9\,\% (2025) frente a +7,5\,\% del total.
\end{itemize}

\sub{II.3. Retos\label{rep:II.3}}
\begin{itemize}
\item Trampa de la renta media (\autor{Gill}-\autor{Kharas}, 2007; BM, 2024); fragmentación; deuda y ciclo financiero \vertema{3.B.29}; materias primas y clima \vertema{3.B.31}.
\end{itemize}

% ==========================================================================
\bloque[8,5]{III. Perspectivas económicas mundiales\label{rep:III}}
\begin{ideaclave}
Resiste una sucesión de choques mejor de lo esperado, pero crece menos que antes; todos los organismos prevén $\approx$3\,\% (PPA) en 2026 y algo más en 2027, con riesgos a la baja y poco margen fiscal; a largo plazo el centro de gravedad se desplaza hacia Asia.
\end{ideaclave}

\sub{III.1. Coyuntura: resistente bajo presión\label{rep:III.1}}
\begin{itemize}
\item WEO oct. 2026, \emph{Resilience Under Strain, Urgent Choices}: 2 fuerzas, \emph{geopolítica} (Ormuz, $\approx$20 mb/d; crudo $\approx$96\,\$, +41\,\%; reapertura supuesta en mar. 2027) y \emph{tecnología} (inversión en IA, choque positivo de demanda): perspectiva «bifurcada».
\item Mundo \dato{3,5 / 3,1 / 3,4\,\%}{2025-2027}, bajo la media de 2000-2019 (3,7); avanzadas 2,1 (EE.\,UU. 2,3; zona del euro 1,1; Japón 0,7; \textbf{España 2,5}); EMED 4,5 $\to$ 3,7 (Oriente Medio y Asia Central $-0{,}3$; China 4,5; India 7).
\item Inflación 4,1 $\to$ 4,7 $\to$ 4,1\,\%. Comercio de mercancías (volumen) \dato{4,2\,\% (2025), 3,9\,\% (2026p)}{OMC}; bienes de IA, 14,8\,\% del comercio; arancel de EE.\,UU. 2,6\,\% $\to$ $\approx$7\,\%.
\end{itemize}
\graficorepaso[0.92\linewidth]{crecimiento-mundial}{Crecimiento del PIB real, 2000-2031: mundo, avanzadas y emergentes}
\graficorepaso[0.95\linewidth]{crecimiento-regiones}{Crecimiento 2025-2027 por regiones (FMI, oct. 2026)}

\sub{III.2. Previsiones de los organismos\label{rep:III.2}}
\begin{center}\footnotesize
\begin{tabular}{@{}lcccc@{}}
\toprule
2026 / 2027 & Mundo & Z. euro & China & Comercio\\
\midrule
FMI (oct.) & 3,1 / 3,4 & 1,1 / 1,2 & 4,5 / 4,0 & 4,6 / 4,7\\
OCDE (sept.) & 2,9 / 3,0 & 1,0 / 1,0 & 4,5 / 4,2 & --\\
Comisión (nov.) & 3,0 / 3,3 & 1,2 / 1,2 & 4,5 / 4,4 & 4,0 / 4,3\\
BM (jun.)$^{*}$ & 2,5 / 2,8 & 0,8 / 1,3 & 4,2 / 4,3 & 2,9 / 3,3\\
OMC (oct.)$^{*}$ & 2,6 / 2,9 & -- & -- & 3,9 / 4,1\\
\bottomrule
\end{tabular}\\
{\scriptsize $^{*}$ PIB a tipos de mercado (BM en PPA: 2,9 / 3,4); OMC: mercancías.}
\end{center}
\begin{itemize}
\item \emph{Coinciden}: $\approx$3\,\% en PPA (2,5-2,6 a tipos de mercado) en 2026, algo más en 2027; EE.\,UU. $\approx$2,2; zona del euro $\approx$1; China 4-4,5. \emph{Difieren}: en el comercio (BM, de junio, antes del auge de los bienes de la IA). Funcas (sept.): 15-16 de 19 servicios ven desfavorable el entorno internacional.
\end{itemize}

\sub{III.3. Riesgos y recomendaciones\label{rep:III.3}}
\begin{itemize}
\item \emph{A la baja} (todos los organismos): guerra más larga (BM: 1,3\,\% en 2026 con tensión financiera); corrección de la IA (valoraciones, apalancamiento, «financiación circular», deuda); tensiones comerciales; clima; tipos soberanos (deuda avanzadas 108 $\to$ 115\,\% del PIB, EMED 74 $\to$ 84\,\%). \emph{Al alza}: productividad de la IA (+0,1-0,8 pp); paradoja de \autor{Solow}; destrucción creativa (\autor{Aghion}, \autor{Howitt}, \autor{Mokyr}, Nobel 2025).
\item \emph{Recomendaciones}: reconstruir el margen fiscal; ayudas energéticas focalizadas (las de precios cuestan 3-22 veces más); monetaria prudente e independiente; reequilibrio (consumo en China, menos déficit en EE.\,UU.), no aranceles; reformas y cooperación.
\end{itemize}

\sub{III.4. Medio y largo plazo\label{rep:III.4}}
\begin{itemize}
\item Medio plazo \dato{3,2\,\% (2028-2031)}{FMI}: demografía, productividad, fragmentación (NMF 80 $\to$ 72\,\%), transición energética. BM: «casi una década de convergencia perdida» sin China ni India.
\end{itemize}
\textbf{Centro de gravedad económico} (\autor{Quah}, 2011):
\begin{itemize}
\item \emph{a) Idea}: un punto del mapa que resume dónde está la actividad.
\item \emph{b) i) Supuestos}: Tierra esfera; PIB (PPA) en el centroide de cada país; centro de masas proyectado. \emph{ii) Desarrollo}: $\vect{u}_i$ unitario por país y
\begin{formula}
\centering \funcionclave{\vect{c}=\frac{\sum_i Y_i\,\vect{u}_i}{\sum_i Y_i}}
\end{formula}
\item[] proyectado en $\vect{c}/\lVert\vect{c}\rVert$ ($\approx$2.700 km bajo la superficie). \emph{iii) Implicaciones}: se acerca a quien crece más; se hunde si la actividad se dispersa.
\item \emph{c) Extensiones}: 700 localizaciones, tipos de mercado; población, emisiones.
\item \emph{d) Evidencia}: Atlántico Norte (1980) $\to$ mar de Noruega (2000) $\to$ Siberia occidental (2025): $\approx$3.800 km al este y, desde 2000, al sur. \autor{Quah}: Atlántico (1980), este de Esmirna (2008), entre India y China (2050).
\item \emph{e) Valoración}: intuitivo; pero la latitud dice poco, depende de la medida del PIB y no explica (lo hace la convergencia).
\end{itemize}
\graficorepaso[0.95\linewidth]{centro-gravedad}{Centro de gravedad económico, 1980-2031 (PIB en PPA)}

\sub{III.5. Balance: DAFO\label{rep:III.5}}
\begin{center}\footnotesize
\begin{tabular}{@{}p{0.47\linewidth}p{0.47\linewidth}@{}}
\cellcolor{tceesalvia}\textcolor{tceeverde}{\textbf{Fortalezas}}: resiste ($\approx$3\,\%); comercio que se reordena; inversión en IA; Asia y los emergentes tiran. &
\cellcolor{tceesalmon}\textcolor{tceerojo}{\textbf{Debilidades}}: 3,2\,\% a medio plazo, productividad floja; deuda y poco margen fiscal; desequilibrios; convergencia frenada fuera de Asia.\\
\cellcolor{tceeazulpalido}\textcolor{tceeazul}{\textbf{Oportunidades}}: productividad de la IA (+0,1-0,8 pp); transición energética; nuevos acuerdos, Sur-Sur, conectores; demografía de India y África. &
\cellcolor{tceegrisclaro}\textcolor{tceedorado}{\textbf{Amenazas}}: guerra y petróleo; corrección de la IA; aranceles y bloques; deuda soberana, dólar y clima.\\
\end{tabular}
\end{center}

% ==========================================================================
\bloque[1]{Conclusión}
\begin{itemize}
\item \textbf{Recapitulación}: «En conclusión,» comercio de manufacturas, servicios digitales y bienes de la IA en 3 polos (gravedad); finanzas entre ricos al ritmo del ciclo global, en dólares; emergentes que convergen; resiste pero crece menos (3,1\,\% FMI, algo menos los demás; 3,2\,\% a medio plazo) y el centro de gravedad va hacia el este.
\item \textbf{Relevancia}: para España, diversificar destinos \vertema{4.A.24} \vertema{4.A.25}.
\item \textbf{Extensiones}: \vertema{3.B.20} \vertema{3.B.21} \vertema{3.B.18} \vertema{3.B.10} \vertema{3.B.33} \vertema{3.B.29}.
\item \textbf{Opinión}: el comercio, más resistente que sus reglas: se reordena antes que desaparecer.
\item \textbf{Idea final}: «Dios aprieta, pero no ahoga»: los choques aprietan, pero la economía sigue creciendo, empujada desde Asia; la duda es si el resto podrá seguirla en un mundo de bloques.
\end{itemize}

% ==========================================================================
\bloque{Autores y fechas}
\hito{1962}{\autor{Tinbergen}: ecuación de gravedad}
\hito{1975}{\autor{Grubel}-\autor{Lloyd}: comercio intraindustrial}
\hito{1979}{\autor{Anderson}: fundamentos de la gravedad}
\hito{1980}{\autor{Dunning}: paradigma OLI}
\hito{1990}{\autor{Lucas}: ¿por qué el capital no va a los pobres?}
\hito{1992}{\autor{Barro}-\autor{Sala-i-Martín}; \autor{Mankiw}-\autor{Romer}-\autor{Weil}}
\hito{1995}{\autor{McCallum}: efecto frontera}
\hito{2001}{\autor{O'Neill}: los BRIC}
\hito{2003}{\autor{Anderson}-\autor{van Wincoop}: resistencia multilateral}
\hito{2007}{\autor{Gill}-\autor{Kharas}: trampa de la renta media}
\hito{2008}{\autor{Disdier}-\autor{Head}: la distancia persiste}
\hito{2011}{\autor{Quah}: centro de gravedad}
\hito{2013}{\autor{Rey}: ciclo financiero global}
\hito{2020}{\autor{Miranda-Agrippino}-\autor{Rey}}
\hito{2025}{\autor{Gopinath} et al.: economías conectoras}
\hito{2026}{Ormuz; WEO \emph{Resilience Under Strain}; BRICS en Nueva Delhi}

% ==========================================================================
\bloque{Datos clave}
\begin{itemize}
\item PIB mundial: \dato{3,5 / 3,1 / 3,4\,\%}{FMI, 2025-2027}; medio plazo \dato{3,2\,\%}{FMI}.
\item Previsiones 2026: FMI 3,1; Comisión 3,0; OCDE 2,9 (PPA); BM \dato{2,5\,\%}{tipos de mercado}.
\item Inflación mundial: \dato{4,7\,\% (2026)}{FMI}; crudo \dato{+41\,\%}{FMI}.
\item PIB en PPA: China \dato{19,6\,\%}{FMI, 2025}; EE.\,UU. 14,6\,\%; avanzadas 39,5\,\%.
\item Comercio: mercancías \dato{26,3 billones \$}{OMC, 2025}; servicios 9,6; volumen \dato{4,2 / 3,9\,\%}{OMC, 2025-2026p}.
\item Manufacturas \dato{69,7\,\%}{BM}; bienes de IA \dato{14,8\,\%}{OMC, 1.\textsuperscript{er} sem. 2026}.
\item IED \dato{1,6 billones \$}{UNCTAD, 2025}; banca \dato{49,5 billones}{BPI, 2026T2}.
\item Saldos corrientes \dato{3,8\,\% PIB mundial}{FMI, ESR}.
\item Dólar \dato{56,7\,\%}{reservas, COFER}; migrantes \dato{304 millones}{ONU, 2024}.
\end{itemize}

% ==========================================================================
\bloque{Preguntas del tribunal}
\begin{enumerate}
\item \textbf{¿Por qué el FMI y el BM dan crecimientos distintos?} PPA frente a tipos de mercado: con pesos de mercado, el FMI también da 2,5\,\%.
\item \textbf{¿Por qué resiste la economía mundial?} Boom de la IA, Asia, comercio desviado más que reducido.
\item \textbf{¿Qué mide el centro de gravedad?} La media del PIB sobre la esfera; su dirección, no su latitud.
\item \textbf{¿Por qué importa la distancia si el transporte es barato?} Información, idioma, confianza, instituciones (\autor{Disdier}-\autor{Head}).
\item \textbf{¿Qué es la resistencia multilateral?} El coste de comerciar con $j$ en relación con todos los demás.
\item \textbf{¿Por qué el capital va de pobres a ricos?} Productividad, riesgo, instituciones; activos seguros en EE.\,UU. (\autor{Lucas}; \autor{Gourinchas}-\autor{Jeanne}).
\item \textbf{¿Trilema o dilema?} Con capital libre, la Fed marca el ciclo aunque se flote (\autor{Rey}).
\item \textbf{¿Qué es un emergente?} No hay definición oficial: FMI, BM, índices, BRICS.
\item \textbf{¿Convergen los emergentes?} Absoluta desde 2000, lenta; condicional sí; trampa de la renta media.
\item \textbf{¿Qué es una economía conectora?} Vietnam, México: unen bloques; resiliencia, no diversificación.
\end{enumerate}

\begin{trampa}
\begin{itemize}
\item China supera a EE.\,UU. en PPA (2016), no a tipos de mercado (16,6 frente a 26\,\%).
\item EMED superan a las avanzadas en PPA en 2009.
\item Volumen de mercancías 2025: 4,2\,\% (revisado), no 4,6\,\%.
\item IED 2025: 1,6 billones (UNCTAD) frente a 1,8 (BM, balanza de pagos).
\item Gravedad: elasticidad a la distancia $\approx-1$; a la renta $\approx1$.
\end{itemize}
\end{trampa}

\end{columnas}
\end{document}
